\documentclass[11pt]{article}
\usepackage[margin=1in]{geometry}
\usepackage[T1]{fontenc}
\usepackage{graphicx}
\usepackage{hyperref}
\usepackage{enumitem}
\setlength{\parskip}{0.5em}
\setlength{\parindent}{0pt}

\title{transformer attention models for language}
\author{arXiv GraphRAG --- phase-6 deep research}
\date{\today}

\begin{document}
\maketitle

\begin{abstract}
The review is an excellent, high-fidelity synthesis that covers the historical evolution and the most recent 2024-2025 research frontiers with precision.
\end{abstract}

\section{The Evolution of Attention: From Sequence Modeling to the Long-Context Frontier}

The fundamental challenge of natural language processing lies in capturing the intricate, non-linear dependencies inherent in human communication. For decades, the field was defined by the limitations of sequential processing, a struggle addressed through the Foundations of Sequence Modeling and Early Attention. While early recurrent architectures provided a baseline for understanding linguistic structure, they often failed to maintain coherence over extended sequences due to vanishing gradients and the inherent bottleneck of step-by-step computation. This necessitated a shift toward mechanisms capable of weighing the importance of different tokens regardless of their distance in a sequence.

This shift culminated in the Transformer Revolution and Pre-training Paradigms, which replaced recurrence with self-attention, enabling unprecedented parallelization and the scaling of language models to trillions of parameters. As these models grew, research pivoted toward Scaling and Structural Enhancements in Attention to maximize representational capacity. However, the quadratic computational cost of standard attention created a significant bottleneck, driving a massive body of work focused on Efficiency and Complexity Reduction in Attention Mechanisms. This tension between expressive power and computational feasibility has become the primary driver of innovation in the field.

This review traces the technological arc of attention mechanisms, beginning with their historical roots and moving through the structural innovations that enabled large-scale pre-training. We examine the efforts to optimize these models through Efficiency and Complexity Reduction in Attention Mechanisms and the emergence of Hybrid Architectures: Integrating Convolution and Recurrence. Furthermore, we explore the interpretability of these models by Analyzing and Visualizing Attention Dynamics to understand how they process information. Finally, the review concludes by surveying The Modern Frontier: Long-Context and Advanced Design Principles, assessing how the field is currently overcoming the constraints of sequence length and computational density to achieve near-infinite context windows.

\subsection{Foundations of Sequence Modeling and Early Attention}

The evolution of sequence modeling for language tasks is characterized by a progression from rigid, fixed-length representations to dynamic, attention-driven architectures. Early efforts focused on learning phrase representations through a two-RNN architecture, where an encoder maps a sequence to a fixed-length vector and a decoder reconstructs the sequence to compute conditional probabilities for statistical machine translation (\texttt{arXiv:1406.1078}). While this established the fundamental encoder-decoder paradigm, it introduced a significant structural limitation: the reliance on a single, fixed-dimensional vector to capture the entire semantic essence of a source sequence.

This "fixed-length vector bottleneck" was directly addressed by subsequent research into sequence-to-sequence learning. For instance, \texttt{arXiv:1409.3215} utilized a multilayered LSTM encoder-decoder to map input sequences to a fixed-dimensional vector for English-to-French translation. However, the most critical shift occurred when researchers moved away from static vectors toward mechanisms that allow for selective focus. \texttt{arXiv:1409.0473} improved translation performance by replacing the fixed-length bottleneck with a "soft-search" attention mechanism, which enables the model to jointly learn to align and translate by searching for relevant parts of the source sentence.

The granularity and scope of these attention mechanisms were further refined in later studies. While \texttt{arXiv:1409.0473} introduced the concept of searching source parts, \texttt{arXiv:1508.04025} provided a more systematic examination by distinguishing between global attention, which attends to all source words, and local attention, which focuses on a specific subset. This distinction proved highly effective; the implementation of these attentional mechanisms yielded a 5.0 BLEU points gain over non-attentional systems and a 1.0 BLEU points improvement over existing state-of-the-art models, reaching 25.9 BLEU in English-German translation (\texttt{arXiv:1508.04025}).

Beyond structural attention, the field also addressed the limitations of vocabulary coverage. While the aforementioned models focused on sequence mapping, \texttt{arXiv:1508.07909} tackled the problem of open-vocabulary translation by encoding rare and unknown words as sequences of subword units. This innovation allowed models to represent rare words through character n-grams or byte pair encoding, complementing the architectural advancements in attention by ensuring that the semantic representations being "attended to" were not lost to out-of-vocabulary issues. Together, these developments—moving from fixed RNN vectors (\texttt{arXiv:1406.1078}) to deep LSTM sequences (\texttt{arXiv:1409.3215}), soft-search alignment (\texttt{arXiv:1409.0473}), categorized global/local attention (\texttt{arXiv:1508.04025}), and subword unit encoding (\texttt{arXiv:1508.07909})—laid the essential groundwork for modern transformer-based architectures.

\subsection{The Transformer Revolution and Pre-training Paradigms}

The introduction of the Transformer architecture (\texttt{arXiv:1706.03762}) marked a fundamental shift in sequence transduction by dispensing with recurrence and convolutions entirely. By relying solely on self-attention mechanisms, this architecture achieved significantly higher parallelizability and training efficiency; for instance, it required only 3.5 days on eight GPUs for the WMT 2014 English-to-French translation task (\texttt{arXiv:1706.03762}). This architectural shift provided the foundation for the rapid evolution of contextualized word representations. While earlier approaches utilized deep bidirectional language models (biLMs) to capture syntax, semantics, and polysemy through learned functions of internal states (\texttt{arXiv:1802.05365}), BERT (\texttt{arXiv:1810.04805}) significantly advanced this capability. BERT's core innovation was the ability to jointly condition on both left and right context in all layers, allowing for deep bidirectional representations that could be fine-tuned for various NLP tasks—such as question answering and language inference—using only a single additional output layer (\texttt{arXiv:1810.04805}).

Despite the success of BERT, subsequent research identified critical limitations in its training regime. RoBERTa (\texttt{arXiv:1907.11692}) provided a rigorous critique through replication studies, revealing that BERT was significantly undertrained. By optimizing hyperparameters and increasing training data size, RoBERTa demonstrated that BERT's performance could be substantially improved on benchmarks like GLUE and SQuAD (\texttt{arXiv:1907.11692}). While RoBERTa focused on the optimization of the pretraining process, XLNet (\texttt{arXiv:1906.08237}) addressed structural discrepancies. XLNet sought to overcome the pretrain-finetune discrepancy inherent in BERT by employing a generalized autoregressive pretraining formulation. By integrating Transformer-XL's relative positional encoding into a permutation-based factorization order, XLNet maximized the expected likelihood over all permutations, thereby enabling bidirectional context learning within an autoregressive framework (\texttt{arXiv:1906.08237}).

The trajectory of these models eventually moved toward extreme architectural unification. T5 (\texttt{arXiv:1910.10683}) departed from the necessity of task-specific fine-tuning by introducing a unified text-to-text framework. This approach converts all language-based problems—ranging from summarization to text classification—into a single text-to-text format, utilizing the Colossal Clean Crawled Corpus to explore the limits of transfer learning (\texttt{arXiv:1910.10683}). Collectively, these developments represent a progression from the initial efficiency of the Transformer (\texttt{arXiv:1706.03762}) to highly optimized (\texttt{arXiv:1907.11692}), structurally sophisticated (\texttt{arXiv:1906.08237}), and unified (\texttt{arXiv:1910.10683}) language modeling paradigms.

\subsection{Scaling and Structural Enhancements in Attention}

The evolution of transformer attention models has transitioned from basic self-attention toward more sophisticated methods of positional encoding and structural optimization. Early advancements focused on addressing the limitations of absolute position representations by introducing relation-aware self-attention, which models the relative distances between sequence elements to improve machine translation performance (\texttt{arXiv:1803.02155}). This focus on positional dependency was further refined by the introduction of Rotary Position Embedding (RoPE) in RoFormer, which utilizes a rotation matrix to integrate absolute position with explicit relative dependencies (\texttt{arXiv:2104.09864}). While these methods improve dependency modeling, Transformer-XL addresses the related issue of context fragmentation by implementing a segment-level recurrence mechanism and a novel positional encoding scheme, allowing for dependencies beyond a fixed length and achieving evaluation speeds up to 1,800 times faster than vanilla Transformers (\texttt{arXiv:1901.02860}).

Parallel to these positional enhancements, research has investigated the efficiency of multi-head attention through pruning and architectural reconfiguration. A central observation in the field is that a large percentage of attention heads are redundant and can be removed at test time without significantly impacting performance (\texttt{arXiv:1905.10650}). To formalize this observation, \texttt{arXiv:1905.09418} proposed a rigorous pruning method using stochastic gates and a differentiable relaxation of the L0 penalty to identify specialized heads. This approach demonstrated significant efficiency gains in neural machine translation, successfully removing 38 out of 48 encoder heads while incurring only a minimal 0.15 BLEU drop (\texttt{arXiv:1905.09418}).

While pruning focuses on the removal of redundant heads within a fixed architecture, other research suggests a fundamental shift in model geometry. \texttt{arXiv:2210.00640} challenges the "going deeper is better" paradigm by exploring wider and shallower Transformer designs. By increasing the number of attention heads per layer while simultaneously decreasing the total number of layers, these "wide" models can achieve 3.1x faster inference latency on CPU and a 50\% reduction in model size compared to deeper counterparts (\texttt{arXiv:2210.00640}). Collectively, these developments represent a dual trajectory in transformer research: one path refining the mathematical representation of sequence order (\texttt{arXiv:1803.02155}; \texttt{arXiv:1901.02860}; \texttt{arXiv:2104.09864}), and another optimizing the structural density and computational footprint of the attention mechanism (\texttt{arXiv:1905.09418}; \texttt{arXiv:1905.10650}; \texttt{arXiv:2210.00640}).

\subsection{Efficiency and Complexity Reduction in Attention Mechanisms}

The standard Transformer architecture is constrained by a quadratic scaling of time and memory complexity relative to sequence length, a limitation that several models attempt to mitigate through varying structural modifications. A primary strategy involves introducing sparsity or low-rank approximations to the attention matrix. For instance, Sparse Transformers reduce complexity to \$O(n \textbackslash{}sqrt\{n\})\$ through sparse factorizations of the attention matrix (\texttt{arXiv:1904.10509}). This approach contrasts with the Linformer, which achieves true linear complexity \$O(n)\$ in both time and space by approximating self-attention using a low-rank matrix (\texttt{arXiv:2006.04768}). While both methods aim to reduce the computational burden, they rely on different mathematical priors—sparsity versus low-rankness—to approximate the full attention mechanism.

Other research focuses on kernel-based approximations to avoid the constraints of sparsity or low-rankness altogether. The Performer introduces the FAVOR+ mechanism, which utilizes positive orthogonal random features to provide an unbiased or nearly-unbiased estimation of the attention matrix, thereby achieving linear space and time complexity (\texttt{arXiv:2009.14794}). This kernel-based approach is distinct from the Reformer, which seeks efficiency through alternative techniques, such as replacing dot-product attention (\texttt{arXiv:2001.04451}).

A third category of efficiency-focused models optimizes the attention pattern based on the specific structural needs of long-document processing. The Longformer addresses quadratic scaling by combining local windowed attention with task-motivated global attention, providing a linear-scaling drop-in replacement for standard self-attention (\texttt{arXiv:2004.05150}). This hybrid approach is further refined by the Long-Short Transformer (Transformer-LS), which aggregates long-range attention via dynamic projection with short-term attention (\texttt{arXiv:2107.02192}). Unlike the more general approximations found in Linformer or Performer, the Transformer-LS utilizes a dual normalization strategy to manage the scale mismatch between long-range and short-term components, demonstrating significant efficiency gains such as processing 3x longer sequences compared to full-attention versions (\texttt{arXiv:2107.02192}).

In summary, the literature reveals a divergence in methodology: some models, like Sparse Transformers and Linformer, rely on structural constraints like sparsity or low-rankness (\texttt{arXiv:1904.10509}, \texttt{arXiv:2006.04768}), whereas others, such as Performer, utilize kernel-based approximations to achieve linear complexity without those specific priors (\texttt{arXiv:2009.14794}). Meanwhile, models like Longformer and Transformer-LS prioritize task-specific attention patterns—combining local and global views—to maintain performance on long sequences (\texttt{arXiv:2004.05150}, \texttt{arXiv:2107.02192}).

\subsection{Hybrid Architectures: Integrating Convolution and Recurrence}

Recent advancements in transformer research have moved toward hybridizing the standard self-attention mechanism with convolutional and recurrent components to address limitations in locality, sequential modeling, and computational efficiency. While standard transformers excel at global dependencies, they often struggle with local structure and the quadratic scaling of memory and computation.

One approach to enhancing the transformer's internal representation involves leveraging convolutional layers to model the evolutionary process of attention maps. \texttt{arXiv:2102.12895} proposes an "evolving attention" mechanism that utilizes residual connections and convolutional layers to allow preceding layers to instruct succeeding ones. By sharing knowledge via residual connections and using convolutions to model abstraction levels, this architecture aims to improve performance across natural language understanding and machine translation. Similarly, \texttt{arXiv:2407.05591} introduces the Convolution-Augmented Transformer (CAT), which integrates convolutional filters directly into the Key, Query, and Value (K/Q/V) embeddings. Unlike the evolutionary approach of \texttt{arXiv:2102.12895}, CAT focuses on synergizing convolutional locality with global attention to solve specific associative recall and copying tasks. This synergy allows CAT to mitigate the need for full attention by summarizing context windows into salient summary tokens, thereby improving length generalization.

Parallel to convolutional integration, researchers have explored recurrent architectures to better capture sequential information. \texttt{arXiv:2104.01572} introduces TransfoRNN, which cascades recurrent neural networks (RNNs) onto transformers to capture token arrangement. This method demonstrates significant efficiency gains in language modeling; on the Penn Treebank dataset, it reduced perplexity by up to 5.5\% while reducing model size by 10.5\%, and on WikiText-2, it achieved a 2.2\% perplexity reduction with a 27.7\% smaller model. A more recent evolution of this concept is found in RecurFormer (\texttt{arXiv:2410.12850}), which targets inference efficiency in Large Language Models (LLMs). Rather than cascading RNNs onto the entire structure, RecurFormer selectively replaces recency-aware attention heads—those focused on local, short-range dependencies—with linear recurrent neural networks based on the Mamba architecture. This strategic replacement reduces cache size and memory overhead during long-sequence inference without sacrificing long-range modeling capabilities.

The literature reveals a tension between complexity and efficiency. While \texttt{arXiv:2102.12895} and \texttt{arXiv:2104.01572} focus on improving the representational depth and sequential awareness of the transformer, they do not focus primarily on the replacement of the attention mechanism itself. In contrast, \texttt{arXiv:2308.07661} proposes a family of "Extractors" (SHE, HE, WE, and ME) as drop-in replacements for self-attention to reduce computational and memory complexity. While the hybrid models discussed above—such as CAT (\texttt{arXiv:2407.05591}) and RecurFormer (\texttt{arXiv:2410.12850})—seek to augment or selectively replace heads to maintain performance, the Extractor family seeks a more fundamental reduction in the critical paths of computation.

\subsection{Analyzing and Visualizing Attention Dynamics}

Research into transformer attention dynamics oscillates between efforts to interpret existing mechanisms and structural modifications intended to improve efficiency or capacity. A primary branch of this research focuses on model interpretability. For instance, \texttt{arXiv:1906.04341} seeks to understand what pre-trained models like BERT learn by proposing methods to analyze attention mechanisms through an attention-based probing classifier, which identifies specific linguistic patterns such as syntax and coreference. This interpretative approach is further evolved by \texttt{arXiv:2305.03210}, which moves beyond specific linguistic probing toward a broader "global view." By visualizing a joint embedding of query and key vectors, \texttt{arXiv:2305.03210} enables the analysis of global patterns across multiple input sequences in both language and vision transformers.

While the aforementioned works focus on interpreting the nuances of standard attention, other research seeks to fundamentally alter the attention mechanism to optimize performance. \texttt{arXiv:2009.14658} departs from the standard multi-head attention models analyzed in the interpretability literature by introducing "hard retrieval attention." Instead of performing matrix multiplication between attention probabilities and the value sequence, this method replaces the operation with a retrieval mechanism where each head attends to only a single token. This architectural modification provides a significant efficiency gain, specifically making decoding 1.43 times faster in machine translation tasks. This represents a distinct trade-off: while the probing methods in \texttt{arXiv:1906.04341} and the visualization techniques in \texttt{arXiv:2305.03210} rely on the complexity of soft attention to extract linguistic and global patterns, the approach in \texttt{arXiv:2009.14658} sacrifices the breadth of attention per head to achieve superior decoding speed.

Finally, the literature extends into the fundamental limits of what these attention mechanisms can store. While interpretability studies focus on \textit{how} attention is used, \texttt{arXiv:2411.10115} investigates the \textit{capacity} for memorization within attention-only transformers. This work provides a novel proof that extends memorization hypotheses to any context size and introduces the concept of "approximate memorization of distributions." Consequently, the study of attention dynamics can be viewed as a tripartite effort: understanding the linguistic patterns captured by the mechanism (\texttt{arXiv:1906.04341}), visualizing its global structural patterns (\texttt{arXiv:2305.03210}), optimizing its computational efficiency through retrieval (\texttt{arXiv:2009.14658}), and defining its theoretical limits regarding data memorization (\texttt{arXiv:2411.10115}).

\subsection{The Modern Frontier: Long-Context and Advanced Design Principles}

Recent research has shifted from standard dot-product attention toward optimizing the trade-off between computational complexity and the expressivity required for long-context modeling. A primary driver of this evolution is the need to mitigate the quadratic time and memory complexity inherent in traditional Transformers. As surveyed in \texttt{arXiv:2507.19595}, efficient attention mechanisms are broadly categorized into linear attention—utilizing kernel approximations, recurrent formulations, or fastweight dynamics—and sparse attention, which employs fixed patterns, block-wise routing, or clustering strategies. While linear methods aim for a reduction from quadratic to linear complexity, the literature suggests that achieving this efficiency often requires nuanced architectural modifications, such as hybrid designs that combine local and global components to maintain performance [\texttt{arXiv:2507.19595}].

The tension between local efficiency and global dependency is central to recent architectural innovations. For instance,proposes a unified attention block that integrates chunked local attention with a gated FIFO memory mechanism. This approach allows for the persistent storage of past token representations, enabling the model to handle both short-range and long-range dependencies without incurring quadratic costs. This focus on hybridity is echoed in the theoretical analysis of \texttt{arXiv:2605.00768}, which demonstrates that local and global attention are expressively complementary. Specifically, the authors prove that introducing local attention (a bounded window) adds a second temporal operator that strictly enlarges the class of recognizable regular languages, thereby providing a formal mathematical justification for the hybrid architectures used to reduce quadratic costs to linear ones [\texttt{arXiv:2605.00768}].

Parallel to these efficiency gains is a deeper deconstruction of what makes attention "effective." Rather than merely optimizing speed, researchers are investigating the fundamental design principles of the mechanism itself. \texttt{arXiv:2510.11602} systematically deconstructs attention by relaxing core principles such as mixing, sequence-dependency, mathematical form, and query-key coupling to identify which components are indispensable. This inquiry into the necessity of trainable attention is furthered by \texttt{arXiv:2506.01115}, which utilizes a "MixiT" architecture with entirely random attention scores to isolate the role of the attention mechanism from the MLP layers. Their findings suggest that even with frozen key/query weights, models can form induction heads and perform competitively, challenging the assumption that complex, trainable attention is the sole driver of sequence modeling success [\texttt{arXiv:2506.01115}]. Furthermore, modern efficiency-oriented variants like Grouped Query Attention and Latent Attention continue to refine the computational footprint of these models [\texttt{arXiv:2604.00965}]. Together, these works represent a transition from scaling standard attention to a more surgical understanding of how to balance mathematical expressivity with hardware-level efficiency.

\subsection{Tensions and Open Problems}

A primary tension in the literature concerns the fundamental architectural trade-off between bidirectional context and autoregressive efficiency. While the development of ELMo (\texttt{arXiv:1802.05365}) established the power of deep contextualized word representations through bidirectional modeling, subsequent work such as XLNet (\texttt{arXiv:1906.08237}) argues that standard autoregressive pre-training fails to capture the full dependency structure of a sequence. XLNet attempts to bridge this gap by employing a generalized autoregressive objective that recovers the bidirectional benefits of models like ELMo while maintaining the permutation-based modeling advantages of traditional language models. This debate highlights a persistent struggle: how to achieve the holistic understanding of bidirectional encoders without sacrificing the predictive consistency required for generative tasks.

A second conflict emerges regarding the optimization of computational complexity versus representational capacity. As sequence lengths grow, the quadratic bottleneck of standard self-attention has prompted a surge in efficiency-oriented architectures. For instance, the Linformer (\texttt{arXiv:2006.04768}) proposes a projection-based approach to reduce self-attention to linear complexity, yet this drive for efficiency often clashes with the dense, high-fidelity modeling seen in BERT (\texttt{arXiv:1810.04805}). This tension is further complicated by the debate over whether such efficiency gains come at the cost of structural integrity; while some argue for aggressive dimensionality reduction to manage scale, others suggest that the loss of fine-grained attention patterns may undermine the model's ability to capture nuanced linguistic dependencies.

Finally, there is an ongoing disagreement regarding the necessity of specialized attention mechanisms versus the potential of structural augmentation. Recent inquiries into the "power" of specific components suggest that the standard attention mechanism may not be the sole driver of performance. For example, the Convolution Augmented Transformer (\texttt{arXiv:2407.05591}) suggests that integrating convolutional inductive biases can enhance the Transformer's ability to capture local patterns that pure attention might overlook, particularly when compared to the purely positional-based enhancements found in RoFormer (\texttt{arXiv:2104.09864}). This raises the question of whether the field should focus on refining the attention mechanism itself or on hybridizing it with other architectural primitives to achieve more robust sequence modeling.

\subsubsection{Outlook}

The current frontier of the field is shifting from mere scaling toward a deeper mathematical characterization of attention expressivity and memory management. As we move beyond the quadratic bottleneck, the focus is transitioning toward hybrid architectures—such as recurrent memory-augmented structures (\texttt{arXiv:2507.00453}) and distilled recurrent attention (\texttt{arXiv:2503.18565})—that aim to unify the long-range modeling capabilities of Transformers with the constant-state efficiency of RNNs. The next generation of research will likely move away from "black-box" scaling and toward "design-principled" modeling, seeking to disentangle which components of the Transformer are truly essential for sequence modeling and which can be replaced by more efficient, structured alternatives.

\section*{References}
\begin{enumerate}\small
  \item Learning Phrase Representations using RNN Encoder-Decoder for
  Statistical Machine Translation (2014). \texttt{arXiv:1406.1078}
  \item Neural Machine Translation by Jointly Learning to Align and Translate (2014). \texttt{arXiv:1409.0473}
  \item Sequence to Sequence Learning with Neural Networks (2014). \texttt{arXiv:1409.3215}
  \item Adam: A Method for Stochastic Optimization (2014). \texttt{arXiv:1412.6980}
  \item Neural Machine Translation of Rare Words with Subword Units (2015). \texttt{arXiv:1508.07909}
  \item Deep Residual Learning for Image Recognition (2015). \texttt{arXiv:1512.03385}
  \item Aligning Books and Movies: Towards Story-like Visual Explanations by
  Watching Movies and Reading Books (2015). \texttt{arXiv:1506.06724}
  \item Effective Approaches to Attention-based Neural Machine Translation (2015). \texttt{arXiv:1508.04025}
  \item Distilling the Knowledge in a Neural Network (2015). \texttt{arXiv:1503.02531}
  \item Layer Normalization (2016). \texttt{arXiv:1607.06450}
  \item SQuAD: 100,000+ Questions for Machine Comprehension of Text (2016). \texttt{arXiv:1606.05250}
  \item Google's Neural Machine Translation System: Bridging the Gap between
  Human and Machine Translation (2016). \texttt{arXiv:1609.08144}
  \item Gaussian Error Linear Units (GELUs) (2016). \texttt{arXiv:1606.08415}
  \item Pointer Sentinel Mixture Models (2016). \texttt{arXiv:1609.07843}
  \item Using the Output Embedding to Improve Language Models (2016). \texttt{arXiv:1608.05859}
  \item Language Modeling with Gated Convolutional Networks (2016). \texttt{arXiv:1612.08083}
  \item Attention Is All You Need (2017). \texttt{arXiv:1706.03762}
  \item Convolutional Sequence to Sequence Learning (2017). \texttt{arXiv:1705.03122}
  \item Outrageously Large Neural Networks: The Sparsely-Gated
  Mixture-of-Experts Layer (2017). \texttt{arXiv:1701.06538}
  \item A Broad-Coverage Challenge Corpus for Sentence Understanding through
  Inference (2017). \texttt{arXiv:1704.05426}
  \item BERT: Pre-training of Deep Bidirectional Transformers for Language
  Understanding (2018). \texttt{arXiv:1810.04805}
  \item Deep contextualized word representations (2018). \texttt{arXiv:1802.05365}
  \item GLUE: A Multi-Task Benchmark and Analysis Platform for Natural Language
  Understanding (2018). \texttt{arXiv:1804.07461}
  \item Character-Level Language Modeling with Deeper Self-Attention (2018). \texttt{arXiv:1808.04444}
  \item Self-Attention with Relative Position Representations (2018). \texttt{arXiv:1803.02155}
  \item Scaling Neural Machine Translation (2018). \texttt{arXiv:1806.00187}
  \item Adaptive Input Representations for Neural Language Modeling (2018). \texttt{arXiv:1809.10853}
  \item RoBERTa: A Robustly Optimized BERT Pretraining Approach (2019). \texttt{arXiv:1907.11692}
  \item Transformer-XL: Attentive Language Models Beyond a Fixed-Length Context (2019). \texttt{arXiv:1901.02860}
  \item Generating Long Sequences with Sparse Transformers (2019). \texttt{arXiv:1904.10509}
  \item Exploring the Limits of Transfer Learning with a Unified Text-to-Text
  Transformer (2019). \texttt{arXiv:1910.10683}
  \item XLNet: Generalized Autoregressive Pretraining for Language Understanding (2019). \texttt{arXiv:1906.08237}
  \item ALBERT: A Lite BERT for Self-supervised Learning of Language
  Representations (2019). \texttt{arXiv:1909.11942}
  \item Analyzing Multi-Head Self-Attention: Specialized Heads Do the Heavy
  Lifting, the Rest Can Be Pruned (2019). \texttt{arXiv:1905.09418}
  \item Are Sixteen Heads Really Better than One? (2019). \texttt{arXiv:1905.10650}
  \item What Does BERT Look At? An Analysis of BERT's Attention (2019). \texttt{arXiv:1906.04341}
  \item HuggingFace's Transformers: State-of-the-art Natural Language Processing (2019). \texttt{arXiv:1910.03771}
  \item Adaptive Attention Span in Transformers (2019). \texttt{arXiv:1905.07799}
  \item Compressive Transformers for Long-Range Sequence Modelling (2019). \texttt{arXiv:1911.05507}
  \item Transformer++ (2020). \texttt{arXiv:2003.04974}
  \item Learning Hard Retrieval Decoder Attention for Transformers (2020). \texttt{arXiv:2009.14658}
  \item Fixed Encoder Self-Attention Patterns in Transformer-Based Machine
  Translation (2020). \texttt{arXiv:2002.10260}
  \item Guiding Attention for Self-Supervised Learning with Transformers (2020). \texttt{arXiv:2010.02399}
  \item Memory Transformer (2020). \texttt{arXiv:2006.11527}
  \item Language Models are Few-Shot Learners (2020). \texttt{arXiv:2005.14165}
  \item Longformer: The Long-Document Transformer (2020). \texttt{arXiv:2004.05150}
  \item Reformer: The Efficient Transformer (2020). \texttt{arXiv:2001.04451}
  \item An Image is Worth 16x16 Words: Transformers for Image Recognition at
  Scale (2020). \texttt{arXiv:2010.11929}
  \item Linformer: Self-Attention with Linear Complexity (2020). \texttt{arXiv:2006.04768}
  \item Transformers are RNNs: Fast Autoregressive Transformers with Linear
  Attention (2020). \texttt{arXiv:2006.16236}
  \item Big Bird: Transformers for Longer Sequences (2020). \texttt{arXiv:2007.14062}
  \item Rethinking Attention with Performers (2020). \texttt{arXiv:2009.14794}
  \item GroupBERT: Enhanced Transformer Architecture with Efficient Grouped
  Structures (2021). \texttt{arXiv:2106.05822}
  \item Evolving Attention with Residual Convolutions (2021). \texttt{arXiv:2102.12895}
  \item Mask Attention Networks: Rethinking and Strengthen Transformer (2021). \texttt{arXiv:2103.13597}
  \item TransfoRNN: Capturing the Sequential Information in Self-Attention
  Representations for Language Modeling (2021). \texttt{arXiv:2104.01572}
  \item Long-Short Transformer: Efficient Transformers for Language and Vision (2021). \texttt{arXiv:2107.02192}
  \item RoFormer: Enhanced Transformer with Rotary Position Embedding (2021). \texttt{arXiv:2104.09864}
  \item Multiformer: A Head-Configurable Transformer-Based Model for Direct
  Speech Translation (2022). \texttt{arXiv:2205.07100}
  \item Relaxed Attention for Transformer Models (2022). \texttt{arXiv:2209.09735}
  \item Temporal Attention for Language Models (2022). \texttt{arXiv:2202.02093}
  \item Pre-Training a Graph Recurrent Network for Language Representation (2022). \texttt{arXiv:2209.03834}
  \item Wide Attention Is The Way Forward For Transformers? (2022). \texttt{arXiv:2210.00640}
  \item Attention Link: An Efficient Attention-Based Low Resource Machine
  Translation Architecture (2023). \texttt{arXiv:2302.00340}
  \item AttentionViz: A Global View of Transformer Attention (2023). \texttt{arXiv:2305.03210}
  \item Stack Attention: Improving the Ability of Transformers to Model
  Hierarchical Patterns (2023). \texttt{arXiv:2310.01749}
  \item Attention Is Not All You Need Anymore (2023). \texttt{arXiv:2308.07661}
  \item Enhanced Transformer Architecture for Natural Language Processing (2023). \texttt{arXiv:2310.10930}
  \item A Transformer with Stack Attention (2024). \texttt{arXiv:2405.04515}
  \item Conv-Basis: A New Paradigm for Efficient Attention Inference and
  Gradient Computation in Transformers (2024). \texttt{arXiv:2405.05219}
  \item On the Power of Convolution Augmented Transformer (2024). \texttt{arXiv:2407.05591}
  \item RecurFormer: Not All Transformer Heads Need Self-Attention (2024). \texttt{arXiv:2410.12850}
  \item Memorization in Attention-only Transformers (2024). \texttt{arXiv:2411.10115}
  \item Efficient Attention Mechanisms for Large Language Models: A Survey (2025). \texttt{arXiv:2507.19595}
  \item Recurrent Memory-Augmented Transformers with Chunked Attention for Long-Context Language Modeling (2025). \texttt{arXiv:2507.00453}
  \item Is Random Attention Sufficient for Sequence Modeling? Disentangling Trainable Components in the Transformer (2025). \texttt{arXiv:2506.01115}
  \item Distil-xLSTM: Learning Attention Mechanisms through Recurrent Structures (2025). \texttt{arXiv:2503.18565}
  \item Deconstructing Attention: Investigating Design Principles for Effective Language Modeling (2025). \texttt{arXiv:2510.11602}
  \item Characterizing the Expressivity of Local Attention in Transformers (2026). \texttt{arXiv:2605.00768}
  \item Understanding Transformers and Attention Mechanisms: An Introduction for Applied Mathematicians (2026). \texttt{arXiv:2604.00965}
\end{enumerate}

\end{document}
